\documentclass[10pt,a4paper]{amsart}
\usepackage[margin=19mm]{geometry}
\usepackage{amsmath,amssymb,mathtools}
\usepackage[scr=rsfs,cal=euler]{mathalfa}
\usepackage{xfrac}
\usepackage[unicode,hidelinks,bookmarks=false]{hyperref}
\newcommand{\ie}{i.e.\ }
\newcommand{\cf}{cf.\ }
\newcommand{\resp}{respectively\ }
\newcommand{\Subsection}{\subsection}
\providecommand{\nobreakdash}{\nobreak\hbox{-}}
\setlength{\emergencystretch}{2em}
\multlinegap0pt
\usepackage{mathtools}
\usepackage{amssymb}
\newtheorem{lemma}{Lemma}
\newcommand{\N}{\mathbb N}
\newcommand{\R}{\mathbb R}
\title{What to Expect When You're Expecting}
\author{Mark Whitmeyer}
\date{Source: arXiv:2606.30400v1 (CC BY 4.0)}
\begin{document}
\maketitle

\noindent\emph{Source and licence.} What to Expect When You're Expecting. Version arXiv:2606.30400v1 (CC BY 4.0), distributed by arXiv under the Creative Commons Attribution 4.0 International licence, http://creativecommons.org/licenses/by/4.0/, as recorded in the arXiv licence metadata for this exact version and shown on the abstract page https://arxiv.org/abs/2606.30400v1. Full licence notice: \texttt{licenses/arxiv-2606.30400-CC-BY-4.0.txt}.\par\smallskip

\noindent\textit{Editorial scope.} Counted unit: the article's lemma lem:finex (source0.tex lines 170--175). The article's complete original proof of it is delivered with it (source0.tex lines 177--283, one \texttt{proof} environment of 107 lines). No mathematics is written, repaired or re-derived: the statement and the proof are the article's own lines, in the article's own notation, and no retained line is substituted or re-flowed.  No further source-local context is needed: every cross-reference of every retained line resolves inside the two delivered spans.  Nothing was omitted from any retained span, so the package carries no omission record.  No retained line contains a citation, so the wrapper carries no bibliography and no bibliography entry.  No retained line contains a figure environment, an image-inclusion command or drawing code, so no image is delivered and the package declares no asset dependency.  Editorial formatting only, in addition: an AMS \texttt{amsart} preamble that restates the article's own declarations and macros used by the retained lines, namely the environments \texttt{lemma} on separate counters, together with the standard packages those lines need.  The retained environments print in this document as Lemma 1 in order of appearance, whereas the article prints the same items with the numbering of its own full document.  The complete native source of the article as distributed in the arXiv e-print for this exact version is bundled unchanged as \texttt{sources/204017/source0.tex}.
\begin{lemma}[Finite extension lemma]\label{lem:finex}
There exist signed weights \(\left(W_N\right)_{N\in\mathcal Z_n}\) such that, for every \(k\)-count type \(m\),
\[
q\left(m\right)=\sum_{N\in\mathcal Z_n}W_N\prod_{a\in A}\binom{N_a}{m_a}.
\]
\end{lemma}

\begin{proof}
Throughout this proof, for \(r\ge0\) we regard \(\binom{x}{r}\) as a polynomial in \(x\), with the convention \(\binom{x}{0}=1\).

Let \(\mathcal M_k\) denote the set of all \(k\)-count types over \(A\). Define the linear map \(K\colon\R^{\mathcal Z_n}\to\R^{\mathcal M_k}\) by
\[
\left(KW\right)\left(m\right)=\sum_{N\in\mathcal Z_n}W_N\prod_{a\in A}\binom{N_a}{m_a}.
\]
We need to show that \(q\in\operatorname{im}K\).

By finite-dimensional duality, \(\operatorname{im}K=\left(\ker K^*\right)^\perp\) with respect to the pairings
\[
\langle W,G\rangle\coloneqq\sum_{N\in Z_n}W_NG(N)
\qquad \text{and} \qquad
\langle q,\varphi\rangle\coloneqq\sum_{m\in M_k}q(m)\varphi(m).
\]
Consequently, it suffices to show that \(\langle q,\varphi\rangle=0\) for every \(\varphi\in\mathbb R^{M_k}\) with \(K^*\varphi=0\). For such a \(\varphi\), define
\[
P_\varphi(N)\coloneqq\sum_{m\in M_k}\varphi(m)\prod_{a\in A}\binom{N_a}{m_a}.
\]
The adjoint is given by \((K^*\varphi)(N)=P_\varphi(N)\), so the condition \(K^*\varphi=0\) says precisely that \(P_\varphi(N)=0\) for every \(N\in Z_n\).

We need a small interpolation fact. Let \(L\in\N\) and \(b\in\left\{-2,-1,0,1,2\right\}\). If a polynomial \(P\) of degree at most \(L\) vanishes on all integer points
\[
\left\{N\in\N^5\colon\lvert N\rvert=2L+1,\omega\left(N\right)=b\right\},
\]
then \(P\) vanishes identically on the affine space
\[
\left\{N\in\R^5\colon\lvert N\rvert=2L+1,\omega\left(N\right)=b\right\}.
\]

We prove this interpolation fact by induction on \(L\). For \(L=0\), \(P\) is constant, and the assertion is immediate. Assume \(L\ge1\). Let \(e_a\) denote the unit vector in the \(a\)-coordinate. Consider the directions \(u\coloneqq e_{-2}+e_2-2e_0\) and \(v\coloneqq e_{-1}+e_1-2e_0\). If \(b\ge0\), also use \(w\coloneqq e_{-1}+e_2-e_0-e_1\). If \(b<0\), use instead \(w'\coloneqq e_{-2}+e_1-e_{-1}-e_0\). The three chosen directions lie in the tangent space \(\left\{x\in\R^5\colon\lvert x\rvert=0,\omega\left(x\right)=0\right\}\), which has dimension \(3\). The chosen directions are linearly independent: for \(u,v,w\), the \(-2\)-coordinate forces the coefficient of \(u\) to vanish, then the \(2\)-coordinate forces the coefficient of \(w\) to vanish, and then the \(-1\)-coordinate forces the coefficient of \(v\) to vanish. The case \(u,v,w'\) is analogous, using the \(2\)-, \(-2\)-, and \(-1\)-coordinates. Therefore, the three chosen directions span the tangent space.

Each chosen direction has the form \(d=p-q\), where \(p,q\in\mathbb N^5\), \(|p|=|q|=2\), and \(\omega(p)=\omega(q)=s\). Explicitly,
\[
\begin{array}{c|c|c|c}
d&p&q&s\\
\hline
u=e_{-2}+e_2-2e_0&e_{-2}+e_2&2e_0&0\\
v=e_{-1}+e_1-2e_0&e_{-1}+e_1&2e_0&0\\
w=e_{-1}+e_2-e_0-e_1&e_{-1}+e_2&e_0+e_1&1\\
w'=e_{-2}+e_1-e_{-1}-e_0&e_{-2}+e_1&e_{-1}+e_0&-1
\end{array}
\]
For \(u\) and \(v\), \(s=0\), so \(b-s=b\). If \(b\ge0\), then \(b\in\{0,1,2\}\) and the third direction is \(w\), so \(b-s=b-1\in\{-1,0,1\}\). If \(b<0\), then \(b\in\{-2,-1\}\) and the third direction is \(w'\), so \(b-s=b+1\in\{-1,0\}\). Thus, \(b-s\in\{-2,-1,0,1,2\}\) in every case. In the next paragraph, fix one row of the table and use its \(d,p,q,s\).

Define \(Q\left(M\right)\coloneqq P\left(M+p\right)-P\left(M+q\right)\). Then \(Q\) has degree at most \(L-1\). Indeed, if \(P_L\) is the top homogeneous part of \(P\), then \(P_L\left(M+p\right)\) and \(P_L\left(M+q\right)\) have the same degree-\(L\) part, namely \(P_L\left(M\right)\), so the degree-\(L\) terms cancel.

If \(M\in\N^5\), \(\lvert M\rvert=2L-1\), and \(\omega\left(M\right)=b-s\), then both \(M+p\) and \(M+q\) lie in the integer slice \(\lvert N\rvert=2L+1\), \(\omega\left(N\right)=b\). Hence, \(Q\left(M\right)=0\) on that lower slice,
\[\left\{M\in\N^5\colon\lvert M\rvert=2L-1,\omega\left(M\right)=b-s\right\}.\]
By the induction hypothesis, \(Q\) vanishes identically on the whole real affine space \(\lvert M\rvert=2L-1\), \(\omega\left(M\right)=b-s\). Now if \(N\) lies in the real affine space \(\lvert N\rvert=2L+1\), \(\omega\left(N\right)=b\), then \(M=N-q\) lies in the preceding real affine space, 
\[\left\{M\in\R^5\colon\lvert M\rvert=2L-1,\omega\left(M\right)=b-s\right\}.\] Therefore, \(Q\left(N-q\right)=0\), i.e. \(P\left(N+d\right)=P\left(N\right)\) for all real \(N\) satisfying \(\lvert N\rvert=2L+1\) and \(\omega\left(N\right)=b\).

For a fixed \(N\) in this affine space and one of the chosen directions \(d\), set \(F\left(z\right)\coloneqq P\left(N+zd\right)\). Since \(d\) is tangent to the affine space, \(N+zd\) lies in the same affine space for every \(z\in\R\). Applying the identity \(P\left(M+d\right)=P\left(M\right)\) with \(M=N+zd\), we obtain \(F\left(z+1\right)=F\left(z\right)\) for every \(z\in\R\). Thus, \(F\) is a polynomial with period \(1\), and hence, \(F\) is constant. Therefore, \(P\left(N+zd\right)=P\left(N\right)\) for every \(z\in\R\). So \(P\) is invariant under every real multiple of each chosen direction. To see that this makes \(P\) constant on the affine slice, let \(N,N'\) lie in the slice. Then \(N'-N\) lies in the tangent space. Since the chosen directions span the tangent space, we may write
\[
N'-N=\alpha d_1+\beta d_2+\gamma d_3
\]
for the three chosen directions \(d_1,d_2,d_3\). Successively applying invariance along \(d_1,d_2,d_3\) delivers \(P(N')=P(N)\) and so \(P\) is constant on the affine slice.

The slice contains an integer point: if \(b\ne0\), take the count vector with \(N_b=1\), \(N_0=2L\), and all other counts zero; if \(b=0\), take the count vector with \(N_0=2L+1\) and all other counts zero. Since \(P\) vanishes on all integer points of the slice, the constant is \(0\). This proves the interpolation fact.

Returning to \(P_\varphi\), we have shown that \(P_\varphi\) vanishes identically on \(|N|=n,\omega(N)=0\). Let \(H_n\coloneqq\left\{N\in\mathbb R^5\colon |N|=n\right\}\). The restriction of the linear functional \(\omega\) to the affine space \(H_n\) is nonconstant; for instance,
\[
\omega\left(0,0,n,0,0\right)=0
\qquad \text{and} \qquad
\omega\left(0,0,n-1,1,0\right)=1.
\]
Choose affine coordinates \((z,h)\) on \(H_n\) with first coordinate \(z=\omega(N)\). In these coordinates the restricted polynomial has the form \(\tilde P(z,h)\) and satisfies \(\tilde P(0,h)=0\) for all \(h\). Viewing \(\tilde P\) as a polynomial in \(z\) whose coefficients are polynomials in \(h\), its constant coefficient is therefore, zero. Hence, \(\tilde P(z,h)=z\tilde R(z,h)\) for some polynomial \(\tilde R\). Equivalently, as functions on the hyperplane \(H_n\), there exists a polynomial \(R\) of degree at most \(k-1\) such that \(P_\varphi(N)=\omega(N)R(N)\).

Now consider the two lines \(N^{\left(2\right)}\left(s\right)\coloneqq\left(0,0,n-s,0,s\right)\) and \(N^{\left(1\right)}\left(s\right)\coloneqq\left(0,0,n-s,s,0\right)\). They satisfy \(\omega\left(N^{\left(2\right)}\left(s\right)\right)=2s\) and \(\omega\left(N^{\left(1\right)}\left(s\right)\right)=s\). Thus, \(P_\varphi\left(N^{\left(2\right)}\left(s\right)\right)=2sR\left(N^{\left(2\right)}\left(s\right)\right)\) and \(P_\varphi\left(N^{\left(1\right)}\left(s\right)\right)=sR\left(N^{\left(1\right)}\left(s\right)\right)\).

Let \(m_0\coloneqq\left(0,0,k,0,0\right)\) be the all-zero \(k\)-count type, and set \(f_0\coloneqq\varphi\left(m_0\right)\), \(f_r^{\left(2\right)}\coloneqq\varphi\left(m_r^{\left(2\right)}\right)\), and \(f_r^{\left(1\right)}\coloneqq\varphi\left(m_r^{\left(1\right)}\right)\) for \(1\le r\le k\). On the first line,
\[
P_\varphi\left(N^{\left(2\right)}\left(s\right)\right)=\sum_{r=0}^k\binom{s}{r}\binom{n-s}{k-r}f_r^{\left(2\right)},
\]
where \(f_0^{\left(2\right)}=f_0\). On the second line,
\[
P_\varphi\left(N^{\left(1\right)}\left(s\right)\right)=\sum_{r=0}^k\binom{s}{r}\binom{n-s}{k-r}f_r^{\left(1\right)},
\]
where \(f_0^{\left(1\right)}=f_0\). At \(N^{\left(0\right)}\coloneqq\left(0,0,n,0,0\right)\), all terms in \(P_\varphi\left(N^{\left(0\right)}\right)\) vanish except the all-zero \(k\)-count type, and, therefore, \(P_\varphi\left(N^{\left(0\right)}\right)=\binom{n}{k}f_0\). Since \(N^{\left(0\right)}\in\mathcal Z_n\), this value is \(0\). Hence, \(f_0=0\).

Using
\[
\left.\frac{d}{ds}\binom{s}{r}\right\rvert_{s=0}=\frac{\left(-1\right)^{r-1}}{r},\qquad (r\ge1),
\]
we obtain, for \(a\in\left\{1,2\right\}\),
\[
\left.\frac{d}{ds}P_\varphi\left(N^{\left(a\right)}\left(s\right)\right)\right\rvert_{s=0}=\sum_{r=1}^k\frac{\left(-1\right)^{r-1}}{r}\binom{n}{k-r}f_r^{\left(a\right)}.
\]
By the definition of \(d_r\),
\[
\sum_{r=1}^kd_rf_r^{\left(a\right)}=\frac{1}{\binom{n}{k-1}}\left.\frac{d}{ds}P_\varphi\left(N^{\left(a\right)}\left(s\right)\right)\right\rvert_{s=0}.
\]
Since both lines meet at \(N^{\left(0\right)}\), we get \(\left.\frac{d}{ds}P_\varphi\left(N^{\left(2\right)}\left(s\right)\right)\right\rvert_{s=0}=2R\left(N^{\left(0\right)}\right)\) and \(\left.\frac{d}{ds}P_\varphi\left(N^{\left(1\right)}\left(s\right)\right)\right\rvert_{s=0}=R\left(N^{\left(0\right)}\right)\). Therefore,
\[
\sum_{r=1}^kd_r\varphi\left(m_r^{\left(2\right)}\right)=\frac{2R\left(N^{\left(0\right)}\right)}{\binom{n}{k-1}}
\qquad \text{and} \qquad
\sum_{r=1}^kd_r\varphi\left(m_r^{\left(1\right)}\right)=\frac{R\left(N^{\left(0\right)}\right)}{\binom{n}{k-1}}.
\]
Consequently,
\[
\begin{aligned}
\left\langle q,\varphi\right\rangle&=\sum_{r=1}^kd_r\varphi\left(m_r^{\left(2\right)}\right)-2\sum_{r=1}^kd_r\varphi\left(m_r^{\left(1\right)}\right)\\
&=\frac{2R\left(N^{\left(0\right)}\right)}{\binom{n}{k-1}}-2\frac{R\left(N^{\left(0\right)}\right)}{\binom{n}{k-1}}=0.
\end{aligned}
\]
Thus, \(q\) annihilates \(\ker K^*\), and hence, \(q\in\operatorname{im}K\).
\end{proof}

\end{document}
